\pagestyle{fancy}
\renewcommand{\chaptermark}[1]{\markboth{#1}{#1}}
% Testate dei capitoli. Sono raccolte in un comando perche' lo stile "plain",
% ridefinito in template.tex con \fancyhf{}, con questa versione di fancyhdr
% cancella i campi in modo globale: dopo un \pagestyle{plain} (abstract,
% introduzione) le testate vanno ridefinite, cosa che fa main/1_capitolo_1.tex.
\newcommand{\testatecapitoli}{%
  \fancyhf{}%
  \renewcommand{\headrulewidth}{0.4pt}%
  \renewcommand{\footrulewidth}{0pt}%
  \fancyhead[RE]{\textbf{\leftmark}}%
  \fancyhead[LE]{\textbf{\chaptername\ \thechapter\ }}%
  \fancyhead[LO]{\textbf{\leftmark}}%
  \fancyhead[RO]{\textbf{\chaptername\ \thechapter\ }}%
  \fancyfoot[R]{\thepage}%
}
\testatecapitoli

\usepackage{acronym}
\usepackage{stfloats}
\usepackage{enumitem}
\usepackage{pdfpages}
\usepackage{pifont}
\usepackage{amsmath,amsfonts}
\usepackage{array}
\usepackage{multirow}
\usepackage{longtable}
\usepackage{tabularx}
\usepackage{makecell}
\usepackage{ragged2e}
\usepackage{pdflscape}
% "scoring" non va spezzato: nel titolo dell'appendice usciva "sco-ring". EEFQ e i nomi
% inglesi degli strumenti (Inventory, Questionnaire) uscivano spezzati con le regole
% italiane ("EE-FQ", "Invento-ry"): niente sillabazione.
% \hyphenation nel preambolo finirebbe nella lingua attiva in quel momento (l'inglese):
% \babelhyphenation registra le eccezioni per l'italiano, che e' la lingua del testo.
% babel-italian tratta l'apostrofo come lettera: "dell'EEFQ" e' una parola sola e va
% elencata a parte, altrimenti esce ancora "dell'EE-FQ".
\babelhyphenation[italian]{scoring EEFQ l'EEFQ dell'EEFQ all'EEFQ nell'EEFQ sull'EEFQ Inventory Questionnaire}
\usepackage[normalem]{ulem}   % \uline: sottolineato dell'Appendice 1
\usepackage{textcomp}
\usepackage[per-mode=symbol,retain-explicit-plus,separate-uncertainty=true,detect-all=true]{siunitx}
\usepackage{xcolor} % options x11names,dvipsnames passed in main.tex (xcolor is loaded by template.tex)

\newcommand{\cmark}{\ding{51}} % ✓
\newcommand{\xmark}{\ding{55}} % ✗

% =============================================================================
% Posizionamento delle tabelle.
% Le tabelle sono flottanti: non si spezzano mai fra due pagine e, se non entrano
% nel punto esatto in cui sono citate, scivolano alla pagina vicina mentre il testo
% scorre a riempire lo spazio. \FloatBarrier (inserito nei capitoli) impedisce che
% si allontanino dalla sezione o dal caso che le descrive.
% =============================================================================
\usepackage{placeins}
\renewcommand{\topfraction}{0.9}        % quanto di una pagina puo' essere occupato
\renewcommand{\bottomfraction}{0.8}     % da flottanti in alto e in basso
\renewcommand{\textfraction}{0.07}      % testo minimo richiesto in una pagina mista
\renewcommand{\floatpagefraction}{0.75} % riempimento minimo di una pagina di sole tabelle
\setcounter{topnumber}{3}
\setcounter{bottomnumber}{2}
\setcounter{totalnumber}{5}

% Indice ed elenco delle tabelle: la casella dei numeri di pagina della classe book
% (1,55 em) e' troppo stretta per i numeri a tre cifre in corpo 12 e dava 40 righe
% "Overfull" di 1,2 pt. Casella e margine destro delle voci allargati.
\makeatletter
\renewcommand{\@pnumwidth}{2.2em}
\renewcommand{\@tocrmarg}{3.2em}
\makeatother

% Terzo passaggio di TeX per i capoversi che non trovano una divisione in righe
% accettabile: un po' di spazio elastico in piu' invece di una riga fuori margine
% (caso «Sonuga-Barke [50]» a pagina 36).
\setlength{\emergencystretch}{1.5em}

% Con il divieto di vedove e orfane alcune pagine finiscono una riga prima: non
% vanno stirate per riempire il piede (dava 150 avvisi "Underfull \vbox").
\raggedbottom

% Niente righe isolate (controllo finale, cap_1.md): vietata la prima riga di un
% capoverso sola in fondo alla pagina (orfana) e l'ultima riga sola in cima alla
% pagina seguente (vedova). LaTeX sposta la riga, e la pagina finisce una riga prima.
% babel-italian riporta i due valori a 3000 (=solo scoraggiate) ogni volta che
% seleziona l'italiano, cioe' a inizio documento e dopo l'abstract inglese: per
% questo i valori vanno agganciati a \extrasitalian, che viene eseguito in quei momenti.
\makeatletter
\addto\extrasitalian{\widowpenalty=10000 \clubpenalty=10000 \@clubpenalty=10000 }
\makeatother

% \paragraph come titolo a se' stante e non a capoverso: serve per le intestazioni
% non numerate del capitolo 4 ("Scheda di sintesi", "Prestazione al Baby-FE", ...)
% Titoli di sezione a bandiera e senza sillabazione: a capo solo fra le parole.
\titleformat{\section}{\normalfont\Large\bfseries\raggedright\hyphenpenalty=10000\exhyphenpenalty=10000}{\thesection}{1em}{#1}
\titleformat{\subsection}{\normalfont\large\bfseries\raggedright\hyphenpenalty=10000\exhyphenpenalty=10000}{\thesubsection}{1em}{#1}
\titleformat{\subsubsection}{\normalfont\normalsize\bfseries\raggedright\hyphenpenalty=10000\exhyphenpenalty=10000}{\thesubsubsection}{1em}{#1}
\titleformat{\paragraph}[hang]{\normalfont\bfseries\raggedright\hyphenpenalty=10000\exhyphenpenalty=10000}{}{0pt}{#1}
\titlespacing*{\paragraph}{0pt}{13.2pt}{*0}

\definecolor{mygreen}{HTML}{007480}
\definecolor{myred}{HTML}{6D1A36}

\lstdefinestyle{python}{
    language=Python,
    basicstyle=\ttfamily\footnotesize,
    keywordstyle=\color{blue}\bfseries,
    commentstyle=\color{gray}\itshape,
    stringstyle=\color{orange},
    showstringspaces=false,
    breaklines=true,
    frame=single,
    columns=fullflexible,
    backgroundcolor=\color{white},
    captionpos=b
}

\usepackage{xurl} % consente di spezzare i DOI lunghi in bibliografia
\usepackage{csquotes}
% doi e url restano nel .bib ma non vengono stampati, come nel template di riferimento:
% la bibliografia risulta piu' compatta. Per rimostrarli basta rimettere true.
% Citazioni multiple compatte in una sola parentesi, [2,4,6], con le sequenze
% compresse in intervalli, [1--3]; la bibliografia resta in stile IEEE.
\usepackage[style=ieee,citestyle=numeric-comp,doi=false,isbn=false,url=false,eprint=false]{biblatex}
\renewcommand*{\multicitedelim}{\addcomma}   % nessuno spazio dopo la virgola
\addbibresource{tail/References.bib}

% Stop biblatex from changing the case of titles/booktitles
\DeclareFieldFormat{titlecase}{#1}
\DeclareFieldFormat{sentencecase}{#1}
\DefineBibliographyExtras{italian}{
  \protected\def\mkbibsentencecase#1{#1}
}

% Metadati del PDF (titolo e autrice nelle proprieta' del file): hyperref e' caricato in
% template.tex, qui si completano solo i campi.
\hypersetup{
  pdftitle={Profili di Disregolazione nella Prima Infanzia: Una Prospettiva Transdiagnostica per la Valutazione Precoce},
  pdfauthor={Silvia Occhionero},
  pdfsubject={Tesi di laurea magistrale in Psicologia dello sviluppo tipico e atipico, Università degli Studi di Genova, A.A. 2025/2026},
  pdflang={it-IT}
}

\renewcommand{\thefootnote}{\Roman{footnote}}

\newcommand{\insertblankpages}[1]{
  \clearpage
  \null
  \thispagestyle{empty}
  \newpage
  \ifnum #1>1
    \insertblankpages{\numexpr#1-1\relax}
  \fi
}
